%! Introduction to Polynomial Functions — Unit 2, Lesson 2.0 — Group Activity (Answer Key)
\documentclass[10pt]{article}
\usepackage{precalculus-article}
\usepackage{precalculus-key}   % pulls in -boxes; do NOT also load -boxes
\newcommand{\TallMath}[1]{$\displaystyle #1\rule[-1.4em]{0pt}{3.2em}$}

\begin{document}

\pageheader{Unit 2, Lesson 2.0}{Group Activity --- Answer Key}

\vspace{0.12in}

\begin{scenariobox}[Sorting the Polynomial Family]{plumlight}
\small
Eight expressions are waiting to be sorted. Some belong to the polynomial family
and some are impostors. Every tier below works from this same bank, so keep it
handy.

\smallskip
\renewcommand{\arraystretch}{1.7}
\begin{tabular}{@{}p{0.23\linewidth}p{0.23\linewidth}p{0.23\linewidth}p{0.23\linewidth}@{}}
\textbf{A.} $4x^3 - x + 2$ &
\textbf{B.} $\dfrac{6}{x^2}$ &
\textbf{C.} $x^2 - 9$ &
\textbf{D.} $\sqrt{x} - 3x$ \\
\textbf{E.} $-7$ &
\textbf{F.} $5x^4 - \dfrac{x}{2} + 1$ &
\textbf{G.} $3^x + x^2$ &
\textbf{H.} $\dfrac{x+1}{x-4}$ \\
\end{tabular}
\end{scenariobox}

\vspace{0.1in}

\boxguard[30]
\begin{tcolorbox}[colback=white, colframe=black!40,
    title=\textbf{Tier R --- Remediate},
    fonttitle=\bfseries, arc=2mm, left=3mm, right=3mm, top=2mm, bottom=2mm]

Sort every expression in the bank. If it \textit{is} a polynomial, give its degree;
if it is not, leave the degree column empty. The first row is done for you.

\smallskip
\renewcommand{\arraystretch}{1.5}
\begin{tabular}{|c|c|c||c|c|c|}
\hline
\textbf{Letter} & \textbf{Polynomial?} & \textbf{Degree} &
\textbf{Letter} & \textbf{Polynomial?} & \textbf{Degree} \\
\hline
A & yes & 3 & E & \ans{yes} & \ans{0} \\
\hline
B & \ans{no} & \ans{---} & F & \ans{yes} & \ans{4} \\
\hline
C & \ans{yes} & \ans{2} & G & \ans{no} & \ans{---} \\
\hline
D & \ans{no} & \ans{---} & H & \ans{no} & \ans{---} \\
\hline
\end{tabular}
\smallskip

\begin{enumerate}[itemsep=6pt]
  \item How many of the eight are polynomials? \quad \ans{4}
  \item Which polynomial in the bank has the \textbf{highest} degree? \quad
        Letter \ans{F}
  \item Which polynomial in the bank has degree \textbf{0}? \quad
        Letter \ans{E}
\end{enumerate}
\end{tcolorbox}

\vspace{0.1in}

\boxguard[30]
\begin{tcolorbox}[colback=white, colframe=black!40,
    title=\textbf{Tier A --- Approaching Proficiency},
    fonttitle=\bfseries, arc=2mm, left=3mm, right=3mm, top=2mm, bottom=2mm]

These polynomials arrived with their terms out of order. Put each in standard form,
then read off its parts.

\smallskip
\renewcommand{\arraystretch}{1.7}
\begin{tabular}{|p{0.19\linewidth}|p{0.26\linewidth}|c|c|c|}
\hline
\textbf{As given} & \textbf{Standard form} & \textbf{Degree} &
\textbf{Leading coef.} & \textbf{Constant term} \\
\hline
$7x - 4 + x^3$ & \ans{$x^3 + 7x - 4$} & \ans{3} & \ans{1} & \ans{$-4$} \\
\hline
$2 - 5x^2$ & \ans{$-5x^2 + 2$} & \ans{2} & \ans{$-5$} & \ans{2} \\
\hline
$x - 6x^4 + 3x^2$ & \ans{$-6x^4 + 3x^2 + x$} & \ans{4} & \ans{$-6$} & \ans{0} \\
\hline
\end{tabular}
\smallskip

\begin{enumerate}[itemsep=8pt]
  \item Name each of the three by its degree family (linear, quadratic, cubic,
        quartic).

        First: \ans{cubic} \quad Second: \ans{quadratic} \quad
        Third: \ans{quartic}

  \item The third polynomial has no term written without an $x$. What is its
        constant term, and how do you know?

        \vspace{0.05in}
        \ansline{It is $0$: the missing term is $+\,0$, and $p(0) = 0$ confirms it.}

  \item Two of the three have a \textbf{negative} leading coefficient. Explain how
        you can spot a negative leading coefficient once a polynomial is in
        standard form.

        \vspace{0.05in}
        \ansline{In standard form the highest-power term comes first, so just read its sign.}
\end{enumerate}
\end{tcolorbox}

\vspace{0.1in}

\boxguard[30]
\begin{tcolorbox}[colback=white, colframe=black!40,
    title=\textbf{Tier E --- Extension},
    fonttitle=\bfseries, arc=2mm, left=3mm, right=3mm, top=2mm, bottom=2mm]

\begin{enumerate}[itemsep=8pt]
  \item Four expressions in the bank are impostors. For each, name the rule it
        breaks: \textit{negative or fractional exponent}, \textit{variable in a
        denominator}, or \textit{variable in an exponent}.

        \smallskip
        \renewcommand{\arraystretch}{1.5}
        \begin{tabular}{|c|p{0.30\linewidth}|p{0.42\linewidth}|}
        \hline
        \textbf{Letter} & \textbf{Rewritten} & \textbf{Rule it breaks} \\
        \hline
        B & \ans{$6x^{-2}$} & \ans{negative exponent (also a denominator)} \\
        \hline
        D & \ans{$x^{1/2} - 3x$} & \ans{fractional exponent} \\
        \hline
        G & \ans{$3^x + x^2$} & \ans{variable in an exponent} \\
        \hline
        H & \ans{no whole-number form} & \ans{variable in a denominator} \\
        \hline
        \end{tabular}

  \item Build a polynomial in standard form with degree $3$, leading coefficient
        $-2$, and constant term $5$.

        \vspace{0.06in}
        $p(x) = $ \ans{$-2x^3 + x^2 + 4x + 5$ (one of many)}

        \vspace{0.04in}
        Is there more than one correct answer? Explain.

        \vspace{0.05in}
        \ansline{Yes --- the middle coefficients are unconstrained, so infinitely many work.}

  \item Marisol says $\dfrac{x^3 - 2x}{4}$ is not a polynomial ``because there is a
        fraction in it.'' Rafi says it is. Who is right, and why?

        \vspace{0.05in}
        \ansline{Rafi. Dividing by the number $4$ only scales the coefficients: this is $\frac{1}{4}x^3 - \frac{1}{2}x$, exponents still whole. Only a variable in the denominator breaks the rule.}

  \item Explain why no polynomial can have degree $\tfrac{1}{2}$, using the
        definition rather than an example.

        \vspace{0.05in}
        \ansline{Degree is the largest exponent on $x$, and the definition requires every exponent to be a whole number. Since $\frac{1}{2}$ is not one, no such term is ever allowed.}
\end{enumerate}
\end{tcolorbox}

\end{document}
